\DocumentMetadata{
  pdfversion=2.0,
  pdfstandard=ua-2,
  lang=en-US,
  tagging=on,
  tagging-setup={math/setup=mathml-SE,tabsorder=structure}
}
\documentclass[11pt,letterpaper]{article}
\usepackage[margin=0.68in,includefoot]{geometry}
\usepackage{newcomputermodern}
\usepackage{amsmath,amscd,mathtools}
\usepackage{tikz,adjustbox}
\providecommand{\square}{\mdlgwhtsquare}
\usepackage[hidelinks]{hyperref}
\hypersetup{
  pdftitle={Analysis PhD Exam, January 2003},
  pdfauthor={University of Florida Department of Mathematics},
  pdfsubject={Graduate mathematics examination},
  pdfdisplaydoctitle=true,
  bookmarksopen=true
}
\setcounter{secnumdepth}{0}
\setlength{\parindent}{0pt}
\setlength{\parskip}{0.28em}
\setlength{\emergencystretch}{3em}
\setlength{\leftmargini}{2.15em}
\setlength{\leftmarginii}{2.65em}
\setlength{\leftmarginiii}{3em}
\renewcommand{\labelenumi}{\textbf{\theenumi.}}
\pagestyle{empty}
\raggedbottom
\allowdisplaybreaks
\newenvironment{examproblems}[1]{%
  \begin{enumerate}%
  \setcounter{enumi}{#1}%
  \setlength{\topsep}{0.35em}%
  \setlength{\partopsep}{0pt}%
  \setlength{\itemsep}{0.48em}%
  \setlength{\parsep}{0.18em}%
}{\end{enumerate}}
\newenvironment{examsubparts}{%
  \begin{enumerate}%
  \setlength{\topsep}{0.18em}%
  \setlength{\partopsep}{0pt}%
  \setlength{\itemsep}{0.16em}%
  \setlength{\parsep}{0.1em}%
}{\end{enumerate}}
\newenvironment{examinstructions}{%
  \par\begingroup\itshape
}{%
  \par\endgroup\vspace{0.18em}
}
\makeatletter
\renewcommand\section{\@startsection{section}{1}{0pt}%
  {-1.25ex plus -.25ex minus -.1ex}%
  {0.55ex plus .1ex}%
  {\normalfont\large\bfseries}}
\renewcommand{\@maketitle}{%
  \newpage\null\vskip -1em%
  {\footnotesize\bfseries
    \href{https://www.ufl.edu/}{University of Florida}, \href{https://math.ufl.edu/}{Department of Mathematics}\par}%
  \vskip -0.5em%
  \tagstructbegin{tag=Title}%
    {\LARGE\bfseries\href{https://gma.math.ufl.edu/exams/analysis-phd/}{Analysis PhD Exam}, January 2003\par}%
  \tagstructend%
  \vskip 0.25em%
  \tagmcbegin{artifact}%
    \hrule height 1pt%
  \tagmcend%
  \vskip 1em%
}
\makeatother
\title{Analysis PhD Exam, January 2003}
\author{}
\date{}
\begin{document}
\maketitle
\thispagestyle{empty}
\section{Part I}
\begin{examinstructions}
State the following theorems.
\end{examinstructions}
\begin{examproblems}{0}
\item
Hölder’s inequality.
\item
The Hahn Decomposition Theorem.
\item
The Fubini Theorem for positive measures~\(\mu\) and~\(\nu\).
\item
The Monotone Class Theorem.
\item
The Hahn–Banach Theorem.
\item
The Monotone Convergence Theorem.
\item
The completeness of \(L_E^1(\mu)\).
\item
The theorem of uniform boundedness.
\end{examproblems}
\section{Part II}
\begin{examinstructions}
Prove one of either Theorem 3 or 7 from part I.
\end{examinstructions}
\section{Part III}
\begin{examinstructions}
Solve the following problems. In problems one through four \((X,\Sigma,\mu)\) is a measure space and~\(F\) is a Banach space.
\end{examinstructions}
\begin{examproblems}{0}
\item
Let \(\mathcal{F}\subset\Sigma\) be a~\(\sigma\)-subalgebra. Prove the existence of the conditional expectation \(E(f\mid\mathcal{F})\) for \(f\in L_F^1(\mu)\).
\item
Let~\(\mathcal{R}\) be a ring generating~\(\Sigma\) and suppose \(f\in L_F^1(\mu)\). Prove, if \(\int_A f\,d\mu=0\) for each \(A\in\mathcal{R}\), then \(f=0\), \(\mu\) almost everywhere.
\item
Suppose~\(E\) (as well as~\(F\)) is a Banach space and~\(U\) is a bounded linear operator from~\(E\) to~\(F\). Show, if \(f\in L_E^1(\mu)\), then \(U\circ f\in L_F^1(\mu)\) and 
\[
\int U\circ f\,d\mu=U\int f\,d\mu.
\]
\item
Suppose~\(X\) is a locally compact Hausdorff space. Show, if~\(K\) is a compact set and the restriction of~\(f\) to~\(K\), \(f|_K\), is continuous, then \(f\phi_K\) is a Borel function. Here \(\phi_K\) is the characteristic function of~\(K\).
\item
Prove that \(C([0,1])\) is a closed proper subspace of \(L^\infty([0,1])\) (the \(L^\infty\) space with respect to Lebesgue measure on the interval \([0,1]\)). Explain how to view \(L^1([0,1])\) as a subspace of \(L^\infty([0,1])^*\). Is \(L^1([0,1])\) a proper subspace of \(L^\infty([0,1])^*\)?
\end{examproblems}
\end{document}
